%%=============================================================================
%% CAPITULO 2 -- REFERENCIAL TEORICO
%%=============================================================================

\section{Referencial teórico}

O Governo Digital consolidou-se como paradigma de reforma administrativa
orientada a dados, deslocando o eixo da modernização do mero provimento
eletrônico de serviços para a reconfiguração dos próprios processos decisórios
da burocracia \parencite{diniz2009,bovens2002}. Nessa trajetória, a eficiência
deixa de ser apreendida como simples redução de custos e passa a articular-se
às noções de inovação no setor público e de capacidade de resposta
institucional, em que a tecnologia opera como vetor de transformação
organizacional, e não como camada meramente acessória \parencite{freire2023}.
A introdução de inteligência artificial nesse percurso, contudo, suscita uma
tensão estruturante (entre ganho de eficiência e preservação de valores
públicos como transparência, equidade e controle) que a literatura recente
sobre governança da IA no setor público identifica como problema central do
campo \parencite{filgueiras2022}.

Parte-se da formulação clássica de \citeonline{bovens2002} sobre a transição
da burocracia de balcão (\textit{street-level bureaucracy}), em que o agente
público exerce juízo individualizado caso a caso, para a burocracia de sistema
(\textit{system-level bureaucracy}), em que rotinas decisórias são
progressivamente estruturadas, padronizadas e absorvidas por sistemas de
informação. Mobiliza-se, ainda, a literatura sobre discricionariedade digital
\parencite{busch2018}, que sistematiza os efeitos da automação sobre a margem
de apreciação dos agentes de ponta.

\subsection{Governança algorítmica e proteção de dados}

A governança algorítmica impõe equilibrar a busca por eficiência operacional
com a segurança jurídica, a explicabilidade das decisões, a proteção de dados
pessoais e a supervisão humana ativa (\textit{human-in-the-loop}) sobre os atos
automatizados da Administração Pública \parencite{filgueiras2022}, exigências
que, no Poder Judiciário, a Resolução CNJ nº~615/2025 passou a disciplinar
\parencite{tamer2025}.

No plano da proteção de dados, a Lei Geral de Proteção de Dados Pessoais
assegura ao titular a revisão de decisões automatizadas, nos seguintes termos:

\begin{citacao}
O titular dos dados tem direito a solicitar a revisão de decisões tomadas
unicamente com base em tratamento automatizado de dados pessoais que afetem
seus interesses, incluídas as decisões destinadas a definir o seu perfil
pessoal, profissional, de consumo e de crédito ou os aspectos de sua
personalidade. \parencite[art.~20]{brasil2018lgpd}
\end{citacao}

\subsection{Experiências na Advocacia Pública}

Localmente, o projeto de pesquisa acadêmica OSIRIS, desenvolvido pela
Universidade de Brasília em parceria com a PGDF e o Tribunal de Justiça do
Distrito Federal e dos Territórios, classifica atos processuais da execução
fiscal \parencite{braz2025,bonat2026}.

Faltam, contudo, avaliações de desempenho com método declarado: no OSIRIS, o
ganho de tempo é estimado sem método explícito \parencite{bonat2026}.
